\documentclass[11pt,a4paper]{article}
\usepackage{webclases-lesson}
\lessonSetup{T4}{v0.3}{T4 - Rozamiento y aplicaciones de las fuerzas}
\begin{document}
\shorthandoff{<>}
\color{Ink}
\coverHeader{T4 --- Rozamiento y aplicaciones de las fuerzas}{Contactos, curvas, peraltes y elasticidad.}
Aplicamos Newton a contactos que pueden impedir o permitir el deslizamiento,
a carreteras curvas y a muelles. La pregunta inicial es qu\'e modelo describe
el estado del cuerpo.
\begin{lessonindex}
\indexentry{sec:b1}{1. Antes y despu\'es de deslizar}
\indexentry{sec:b2}{2. Tomar una curva}
\indexentry{sec:b3}{3. Muelles y medidas}
\end{lessonindex}
En los ejemplos mec\'anicos usamos $g=9,81\mpss$. Las estimaciones experimentales de
$g$ se calculan a partir de sus datos.
\precisionnote

\block{1}{Antes y despu\'es de deslizar}{sec:b1}
El rozamiento se opone al deslizamiento relativo, o a su tendencia, en el contacto.
No tiene por qu\'e oponerse a la velocidad del centro de masas en cualquier situaci\'on.
Aqu\'i estudiamos bloques sobre superficies fijas:
\[
|f_s|\le\mu_sN,\qquad |f_k|=\mu_kN.
\]
El est\'atico $f_s$ se ajusta hasta un m\'aximo; el cin\'etico $f_k$ se usa
cuando hay deslizamiento. Los coeficientes son aproximaciones emp\'iricas del contacto.
Primero se calcula $N$; despu\'es la fuerza necesaria para impedir el deslizamiento;
solo entonces se comprueba si cabe dentro del l\'imite est\'atico.

\begin{examplebox}{misma fuerza, dos estados}
\origin{Elaboraci\'on propia, conservando los datos de la versi\'on actual.}
\step{Planteamiento} Bloque de $2,00\kg$ en mesa horizontal,
$\mu_s=0,40$, $\mu_k=0,30$ y fuerza $F=7,00\N$ a la derecha.
Calcula rozamiento y aceleraci\'on desde reposo y si ya desliza a la derecha.
\solutionblock
\step{Modelo} Verticalmente $N=mg$. Horizontalmente, desde reposo se necesita
$f_s=F$ para equilibrar; si desliza usamos $f_k=\mu_kN$.
\step{C\'alculo} $N=19,6\N$, $f_{s,\max}=7,85\N$, $f_k=5,89\N$.
Desde reposo $7,00<7,85$: $f_s=7,00\N$ hacia la izquierda y $a=0$.
Deslizando: $a=(7,00-5,89)/2=+0,557\mpss$.
\step{Interpretaci\'on} El valor de $F$ no decide por s\'i solo el estado.
Usar siempre $f_s=\mu_sN$ predecir\'ia una fuerza neta hacia la izquierda en un
bloque que nadie ha puesto en movimiento.
\begin{center}
\begin{tikzpicture}[x=.48cm,y=.38cm,font=\small]
\draw[axis] (0,0)--(13,0) node[right] {$F$ (N)};
\draw[axis] (0,0)--(0,10) node[above] {$|f|$ (N)};
\draw[BrandBlue,thick] (0,0)--(7.848,7.848);
\draw[GreenLine!70!black,dashed,thick] (7.848,5.886)--(12,5.886);
\draw[guide] (7.848,0)--(7.848,7.848);
\fill[BrandBlue] (7.848,7.848) circle (2pt);
\draw[fill=white,draw=GreenLine!70!black] (7.848,5.886) circle (2pt);
\foreach \x in {4,8,12} \draw (\x,0)--(\x,-.15) node[below] {$\x$};
\foreach \y/\lab in {5.886/{5,89},7.848/{7,85}} \draw (0,\y)--(-.15,\y) node[left] {$\lab$};
\node[above,rotate=38] at (4,4.4) {reposo};
\node[above] at (10,6.1) {desliza};
\end{tikzpicture}
\end{center}
La gr\'afica sigue una fuerza que aumenta lentamente desde reposo.
Al superar $7,85\N$ cambia el r\'egimen: el descenso a $5,89\N$ no representa una
rama inclinada continua. Un bloque ya deslizante puede tener $f_k=5,89\N$
tambi\'en para fuerzas menores.
\end{examplebox}

\heading{Qu\'e pasa si se retira la fuerza}
Si el bloque anterior desliza a $3,00\mps$ y retiramos $F$, queda
$a=-f_k/m=-2,94\mpss$. Se para en $1,02\units{s}$ tras $1,53\units{m}$.
Desde ese instante el rozamiento est\'atico necesario es cero.
Prolongar la aceleraci\'on negativa lo har\'ia arrancar hacia atr\'as sin causa.

\begin{examplebox}{reposo y deslizamiento en el plano}
\origin{Adaptaci\'on del plano inclinado del original T4, con los datos actuales.}
\step{Planteamiento} $m=0,300\kg$, $\theta=25^\circ$, $\mu_s=0,50$ y
$\mu_k=0,35$. Decide si permanece en reposo y calcula $a$ si se lanza cuesta abajo.
\solutionblock
\step{Modelo} Eje $+s$ cuesta abajo. Sin aceleraci\'on perpendicular:
$N=mg\cos\theta$. El peso paralelo apunta cuesta abajo, el rozamiento
hacia arriba en los dos estados que comparamos.
\begin{center}
\begin{tikzpicture}[x=1cm,y=1cm,font=\small]
\draw[fill=Panel,draw=Border] (0,0)--(5.5,0)--(5.5,{5.5*tan(25)})--cycle;
\begin{scope}[rotate=25]
 \draw[fill=BlueSoft,draw=BrandBlue] (2.1,0) rectangle (2.9,.55);
 \coordinate (B) at (2.5,.275);
 \draw[force] (B)--++(0,{1.6*cos(25)}) node[above] {$N$};
 \draw[vectoramber] (B)--++({1.6*sin(25)},0) node[above right] {$f_s$};
 \draw[guide] (B)--++({-1.6*sin(25)},0) node[left] {$mg\sin\theta$};
 \draw[axis] (1.7,-.85)--(.4,-.85) node[left] {$+s$};
\end{scope}
\draw[force] (B)--++(0,-1.6) node[below] {$mg$};
\draw (.8,0) arc[start angle=0,end angle=25,radius=.8];
\node at (1.2,.23) {$25^\circ$};
\end{tikzpicture}
\end{center}
Dibujo del equilibrio: $f_s=mg\sin\theta$. La longitud de cada fuerza sigue
sus componentes; la gu\'ia del peso no es una fuerza adicional.
\step{C\'alculo}
\[
mg\sin25^\circ=1,24\N,\quad N=2,67\N,\quad f_{s,\max}=1,33\N.
\]
Puede quedar en reposo: $f_s=1,24\N<1,33\N$.
Si ya baja, $f_k=\mu_kmg\cos25^\circ=0,934\N$ y
\[
a=g(\sin25^\circ-\mu_k\cos25^\circ)=1,03\mpss
\]
usando trigonometr\'ia sin redondear.
\step{Interpretaci\'on} El l\'imite de reposo exige
$mg\sin\theta\le\mu_smg\cos\theta$, es decir, $\tan\theta\le\mu_s$.
Aqu\'i $\theta_c=\arctan0,50=26,6^\circ$. Es un equilibrio l\'imite ideal.
\variation{Si se lanza cuesta arriba, el rozamiento apunta cuesta abajo:
$a_s=g(\sin25^\circ+\mu_k\cos25^\circ)=7,26\mpss$.
Al parar, vuelve a ser posible el reposo porque $25^\circ<\theta_c$.}
\end{examplebox}

\block{2}{Tomar una curva}{sec:b2}
Para girar a rapidez constante hace falta aceleraci\'on radial $v^2/R$
(T2). Newton exige una resultante hacia el centro, no una fuerza nueva llamada
centr\'ipeta. La normal se orienta seg\'un la superficie, no seg\'un la inclinaci\'on
de un conductor o del veh\'iculo.

\heading{Curva plana}
Sin aerodin\'amica ni aceleraci\'on vertical, $N=mg$. Si no hay derrape,
el rozamiento est\'atico proporciona la fuerza radial:
\[
f_s=\frac{mv^2}{R}\le\mu_smg
\quad\Rightarrow\quad v_{\max}=\sqrt{\mu_sgR}.
\]
No se utiliza $\mu_k$ mientras el contacto se mantenga sin deslizamiento lateral.

\heading{Peralte sin rozamiento}
La normal se inclina porque la calzada est\'a inclinada. Su componente vertical
equilibra el peso y la horizontal sostiene el giro:
\[
N\cos\theta=mg,\qquad N\sin\theta=\frac{mv^2}{R}.
\]
Al dividir, $\tan\theta=v^2/(Rg)$.
La rapidez $v_d=\sqrt{Rg\tan\theta}$ es la de dise\~no \emph{sin rozamiento};
no es una m\'axima universal.

\begin{examplebox}{misma curva, dos modelos}
\origin{Adaptado del original T4, pp.\ 6--7.}
\step{Planteamiento} $R=30,0\units{m}$. Compara curva plana con $\mu_s=0,80$
y peralte a $10^\circ$ sin rozamiento.
\solutionblock
\step{Modelo} Los siguientes son cortes transversales; el centro de la curva
est\'a a la izquierda, y la velocidad es perpendicular a este plano.
\begin{center}
\begin{tikzpicture}[x=1cm,y=1cm,font=\small]
\node at (0,2.4) {plana};
\draw (-1.6,-.4)--(1.6,-.4);
\fill[BrandBlue] (0,0) circle (2.5pt);
\draw[force] (0,0)--(0,1.4) node[above] {$N$};
\draw[force] (0,0)--(0,-1.4) node[below] {$mg$};
\draw[vectoramber] (0,0)--(-1.12,0) node[above left] {$f_s$};
\begin{scope}[xshift=6.5cm]
 \node at (0,2.4) {peralte $10^\circ$};
 \draw (-1.6,{-.4-1.6*tan(10)})--(1.6,{-.4+1.6*tan(10)});
 \fill[BrandBlue] (0,0) circle (2.5pt);
 \draw[force] (0,0)--++(100:{1.4/cos(10)}) node[above] {$N$};
 \draw[force] (0,0)--(0,-1.4) node[below] {$mg$};
 \draw[guide] (0,1.4)--({-1.4*tan(10)},1.4);
 \node[left] at (-.5,.75) {$N\sin\theta$};
\end{scope}
\end{tikzpicture}
\end{center}
\step{C\'alculo} En la plana:
$v_{\max}=\sqrt{0,80(9,81)(30)}=15,3\mps$.
En el peralte:
$v_d=\sqrt{30(9,81)\tan10^\circ}=7,20\mps$.
\step{Interpretaci\'on} En el primer caso limita la adherencia; en el segundo,
la geometr\'ia fija una rapidez compatible sin ayuda del rozamiento.
Con rozamiento, un peralte puede admitir un intervalo de velocidades:
su sentido se decide por la tendencia a deslizar, no se fija siempre cuesta abajo.
\variation{Si en la curva plana $\mu_s$ baja a $0,50$, el l\'imite es
$12,1\mps$. Son resultados de un modelo did\'actico, no velocidades recomendadas.}
\end{examplebox}

\block{3}{Muelles y medidas}{sec:b3}
\heading{Deformaci\'on no es desplazamiento desde el equilibrio}
En el r\'egimen el\'astico lineal:
\[
F_{\mathrm{el}}=-kx.
\]
$x$ mide deformaci\'on respecto a la longitud natural y $k$ se expresa en N/m.
El signo indica fuerza restauradora: se opone a la deformaci\'on.
Para mantener el muelle estirado en equilibrio, la fuerza aplicada opuesta vale $kx$.
No se debe prolongar esta ley fuera del rango lineal del muelle real.

\begin{examplebox}{dos referencias en un muelle vertical}
\origin{Elaboraci\'on propia.}
\step{Planteamiento} $k=100\units{N/m}$ y $m=0,200\kg$.
Calcula el alargamiento de equilibrio y la aceleraci\'on inicial al bajar
la masa $1,00\units{cm}$ m\'as y soltarla.
\solutionblock
\step{Modelo} Con eje positivo hacia abajo, $F_{\mathrm{neta}}=mg-kx$.
En equilibrio, $x_{\mathrm{eq}}=mg/k$. Despu\'es,
$x=x_{\mathrm{eq}}+\xi$, donde $\xi$ se mide desde el equilibrio.
\step{C\'alculo} $x_{\mathrm{eq}}=0,0196\units{m}$.
Al soltar, $x=0,0296\units{m}$ y $\xi=0,0100\units{m}$:
$F_{\mathrm{neta}}=1,96-2,96=-1,00\N$ y $a=-5,00\mpss$.
\begin{center}
\begin{tikzpicture}[x=1cm,y=1cm,font=\small]
\foreach \p/\len/\lab in {0/1.4/natural,4/2.4/equilibrio,8/2.9/al soltar} {
 \draw (\p-.45,0)--(\p+.45,0);
 \draw[thick] (\p,0)--(\p,-\len);
 \node[above] at (\p,0) {\lab};
 \draw[fill=BlueSoft,draw=BrandBlue] (\p-.3,-\len) rectangle ++(.6,-.35);
}
\draw[guide] (-.6,-1.4)--(9,-1.4);
\draw[guide] (3.3,-2.4)--(9,-2.4);
\draw[<->] (9.3,-1.4)--(9.3,-2.9) node[midway,right] {$x=2,96\units{cm}$};
\node[below] at (4,-2.8) {$x=1,96\units{cm}$};
\node[below] at (8,-3.35) {$\xi=1\units{cm}$};
\end{tikzpicture}
\end{center}
Esquema de longitudes: las diferencias verticales siguen aproximadamente la proporci\'on $1,96:1$;
los segmentos representan el muelle, no su forma helicoidal.
\step{Interpretaci\'on} Sustituyendo $x=x_{\mathrm{eq}}+\xi$ resulta
$F_{\mathrm{neta}}=-k\xi$. El peso sigue actuando; queda incorporado al elegir
como referencia el equilibrio. T5 usar\'a la deformaci\'on para calcular energ\'ia el\'astica.
\end{examplebox}

\heading{Una pendiente mide la rigidez}
En equilibrio est\'atico, la gr\'afica de fuerza \emph{aplicada} frente a elongaci\'on
tiene pendiente $+k$; la fuerza \emph{del muelle} tiene pendiente $-k$.
Para el muelle de $100\units{N/m}$:
\begin{center}
\begin{tikzpicture}[x=110cm,y=1cm,font=\small]
\draw[axis] (0,0)--(.05,0) node[right] {$x$ (m)};
\draw[axis] (0,0)--(0,4.5) node[above] {$F_{\mathrm{aplicada}}$ (N)};
\draw[BrandBlue,thick] (0,0)--(.04,4);
\foreach \x/\f in {.01/1,.02/2,.03/3,.04/4} {
 \fill (\x,\f) circle (2pt);
 \draw (\x,0)--(\x,-.08) node[below] {$\x$};
 \draw (0,\f)--(-.0005,\f) node[left] {$\f$};
}
\draw[guide] (.01,1)--(.03,1)--(.03,3);
\end{tikzpicture}
\end{center}
La pendiente es $\Delta F/\Delta x=(3-1)/(0,03-0,01)=100\units{N/m}$.
Estos puntos son calculados del modelo, no medidas realizadas.
Una curvatura sistem\'atica en datos reales advertir\'ia que una sola pendiente ya no basta.

\heading{Leer los datos experimentales del original}
Para un p\'endulo de peque\~na amplitud, el modelo proporcionado es
$T=2\pi\sqrt{L/g}$. Midiendo $L$ desde la suspensi\'on al centro de la masa y
cronometrando varias oscilaciones se obtiene un periodo medio.
Al elevar al cuadrado:
\[
T^2=\frac{4\pi^2}{g}L.
\]
Por eso se representa $T^2$ frente a $L$, no se interpreta la pendiente de $T(L)$
como $4\pi^2/g$. La pendiente publicada en el original,
$4,015\units{s^2/m}$, implica $g=4\pi^2/4,015=9,83\mpss$.
Se usa ese ajuste publicado, no se afirma haberlo reproducido.

Para un muelle, el modelo ideal dado es $T^2=(4\pi^2/k)m$.
El original publica $T^2=(0,0117\units{s^2/g})m+0,0664\units{s^2}$.
La conversi\'on de pendiente da $11,7\units{s^2/kg}$ y
$k=4\pi^2/11,7=3,37\units{N/m}$.
El intercepto no nulo indica que el modelo de partida requiere revisi\'on:
puede haber masa efectiva o sesgos, pero no demuestra una causa \'unica.
Bajo la hip\'otesis adicional $T^2=(4\pi^2/k)(m+m_{\mathrm{ef}})$,
$m_{\mathrm{ef}}=0,0664/11,7=0,00568\kg$.

\begin{exercisebox}{comprobar una tabla y un ajuste}
En la tabla del muelle aparece $T=0,983\units{s}$ junto a
$T^2=0,0966\units{s^2}$. Recalcula el cuadrado con tres cifras significativas.
Explica si la diferencia entre $g=9,83$ y $9,81\mpss$ mide la incertidumbre.
\solutionblock
$0,983^2=0,966289$, luego $T^2=0,966\units{s^2}$:
el valor transcrito tiene un factor diez incorrecto.
La diferencia $0,02\mpss$ es una discrepancia frente a una referencia;
la incertidumbre exige informaci\'on sobre medidas, instrumentos y ajuste.
\end{exercisebox}

\heading{Fuentes y procedencia}
Plano y curvas adaptan el original T4, titulado \emph{Fuerzas centrales};
su contenido real incluye rozamiento y elasticidad. La mesa de $2\kg$, el muelle
vertical y las variaciones son propuestas propias. Los ajustes de p\'endulo y
muelle proceden de sus pp.\ 8--13; no se han realizado experimentos nuevos.
\source{https://fisquiweb.es/Apuntes/apun1B.htm}{FisQuiWeb / IES La Magdalena: inspiraci\'on y datos del original.}
Contraste externo sobre ajuste est\'atico del rozamiento y fuerzas en curvas:
\source{https://openstax.org/books/university-physics-volume-1/pages/6-2-friction}{OpenStax, University Physics 1, 6.2: Friction.}
\source{https://openstax.org/books/university-physics-volume-1/pages/6-3-centripetal-force}{OpenStax, University Physics 1, 6.3: Centripetal Force.}
\end{document}
